\documentclass[11pt,letterpaper]{article}
\usepackage{report-style}

\begin{document}
\reporttitle

\lead{摘要}
本研究围绕大语言模型投资智能体的事件理解问题，探索如何组织可持续修订的金融背景知识，支持相关机制的调用、具体情境的解释与判断的更新。截至2026年11月13日，已建立包含240个概念与机制节点的金融认知图谱，完成两部教材的首轮接入，建设官方与机构材料库及206篇专业分析文章的本地材料库，并完成AI信用融资、韩国股市和美债基差交易三项内容制作试点。两轮探索性对照实验均未观察到使用当前图谱带来的得分增益。阶段工作形成了知识原型、来源处理方法和实验记录；后续将完成试点内容整合，积累具体主体与时间线，并通过更有区分度的对照检验知识积累的作用。

\section{研究问题与总体思路}

\lead{研究目标。}
本项目关注投资智能体从接收事件信息到形成分析判断的过程：识别新增信息，选择相关背景，结合事前预期解释其含义，并在新证据出现后修正判断。金融知识建设是本阶段的研究切入点，目的在于形成可检索、可追溯、可修订的背景体系，研究这种持续积累能否改善事件理解。

\lead{核心问题。}
同一经济变化可能因主体、合同安排、融资条件和时间尺度不同而产生不同影响。一般机制能够提供解释线索，但具体应用仍需辨认当时的约束与信息。本项目据此提出工作假设：将一般知识与具体主体、历史状态及此前判断连接起来，有助于智能体选择相关背景、比较解释并更新认识。该假设的有效性仍需实验检验。

\lead{建设路线。}
研究分四个阶段推进：首先利用模型先验知识构建可讨论的原型，通过样例校准知识组织原则；其次以高质量教材和系统材料修订基础认识；随后通过新闻与专业分析积累具体对象、关系及时间线；最后利用教师课件等材料检验覆盖与解释能力。原型保持暂定性，后续材料可以推动内容和结构的调整。

\lead{阶段进展。}
目前已完成原型及两部教材的首轮接入，进入事件材料的制作与整合阶段。三项新闻试点已形成知识解释和分析稿，尚待整合为正式条目。另已完成两轮探索性对照实验。相较开题计划，本阶段重点推进了知识体系和来源处理方法的实践；系统文献比较与最终研究问题的收敛仍需继续。

\newpage
\section{知识体系设计与基础建设}

\subsection{概念、关系与时间的组织方式}

\lead{概念与具体对象。}
概念页解释对象的含义、构成、制度安排和观测方法，并积累其演变。例如“中央对手方与集中清算”说明清算机构如何介入交易，净额处理覆盖哪些敞口，保证金与违约管理承担哪些职责。独立建页依据其解释职责和后续积累价值，普通术语可以在相关正文中解释。

\lead{机制与具体关系。}
机制页解释变化如何通过主体行为和约束产生结果：谁受到影响，其收益、风险或可行选择如何改变，随后作出什么回应，以及结果依赖哪些条件。事件分析还会揭示特定主体间的融资、交易或持仓联系。这些关系可先在具体情境中展开，经过比较后再判断其可推广范围。

\lead{时间记录与知识继承。}
图谱通过相互链接的Markdown条目组织知识。一般定义和机制在主要位置维护，具体对象可继承通用解释，集中记录自身差异与历史。例如AI相关企业CDS复用CDS的工具基础，重点积累企业融资角色、保护需求和报价分化。时间记录保留信息时点、原始观察、当时判断与后续修正；预期需明确主体、对象、期限及条件。该结构仍处于试点校准阶段。

\subsection{原型与教材接入成果}

原型通过少量样例的反复讨论明确解释深度、对象与机制的分工及选材原则，再扩展为163个概念、71个机制，共234个节点。随后采用Carlin与Soskice的宏观经济学教材，以及Cecchetti与Schoenholtz的货币银行教材开展首轮修订\cite{carlin,cecchetti}。阅读范围分别为导言及第1—18章、实际取得的第6版ISE第1—23章，共组织36个阅读单元。

本轮提出72份候选，筛选保留29份，制作24份全文稿，经整合接入23个页面，其中新增6篇、更新17篇。当前正式图谱包含165个概念、75个机制，共240个节点。

\begin{center}
\small
\begin{tabularx}{\textwidth}{@{}L{0.17\textwidth}Y@{}}
\toprule
成果类型 & 代表性内容 \\
\midrule
新增概念（2篇） & 中央对手方与集中清算；雇主工资制定权。 \\
新增机制（4篇） & 衰退后的就业疤痕；信贷配给；汇率贬值后的贸易收支调整；资源繁荣与非资源出口竞争力。 \\
更新旧页（17篇） & 金融加速器、工资—价格反馈、银行信贷的资本约束、金融周期、资本流动与资本管制等。 \\
\bottomrule
\end{tabularx}
\end{center}

例如，金融加速器的更新区分住房权益提取、信贷条件与待购房者首付储蓄，说明住房价格变化对不同家庭的作用；集中清算条目同时解释双边敞口减少和风险向共同节点集中。采用内容记录版本与页码，后续教材仍可改写已有解释及其组织方式。

\newpage
\section{事件材料建设与内容制作试点}

\subsection{材料基础与制作方法}

\lead{材料库建设。}
官方与机构材料围绕美国通胀、欧洲通胀、中国需求与地产、美国财政融资、日元套息、黄金需求和贸易预期七条线组织，现有71份可读材料，另8份尚未取得正文。专业分析库收录206篇华尔街见闻公开文章，页面日期覆盖2025年1月1日至2026年11月4日，包含770处正文插图，并提供离线阅读、Markdown正文和全文检索。材料取得状态与内容核读、采用状态分别记录。

\lead{从材料形成知识。}
教材以连续章节和完整论证为阅读单位，经候选筛选、全文制作、审阅与返修形成条目；课程幻灯片需追踪跨讲依赖，并结合图注、条件和反例核读。专业分析文章先通读正文与图表，列出重要对象和关系，对照已有图谱，补写必要基础解释，再分析文章试图解释的变化、参与者安排及推断。内容明确后决定新建、更新或复用，允许来源改变原有结构。

\lead{试点组织。}
已完成三篇文章的试作，每篇由一个继承项目研究原则与上下文的子智能体连续处理，再统一审阅。成果共覆盖37项概念清单和18项关系清单，包含已有知识的复用、更新建议及新对象解释；这些数量不代表新增节点。当前已完成试作审阅，正式图谱整合尚在推进。

\subsection{代表性案例及其方法意义}

\lead{AI信用融资：连接工具与具体风险。}
AI CDS文章\cite{aicds}涉及企业CDS、信用指数CDX、供应商融资、代理对冲和信用基差。试作先补齐工具与参与者知识，再解释客户偿付风险如何与供应商订单、债权及担保相连，并比较信用基本面与保护需求对价格变化的不同解释。一般CDS知识与AI企业的具体融资联系分层维护，避免每次事件都重复工具定义。

\lead{韩国股市：区分估值状态与实际承接。}
韩国市场文章\cite{korea}促使试作补写KOSPI、SOX、杠杆ETF及持仓口径，并展开企业回购、外资和境内机构的买盘联系。分析关注估值吸引力、资金约束与实际订单的差别；ETF持仓存量、资金净流量与日度再平衡分别说明。指数比较还需辨认样本与行业暴露，图形共变本身不足以确立因果关系。

\lead{美债基差交易：复用机制并还原交易安排。}
美债文章\cite{basis}已有较多基础知识可复用，增量主要在资管机构、对冲基金与交易商的分工，以及融资、保证金和中介资产负债表约束。试作同时辨认新闻时点与引用研究时期的差异，检查图表和回归结果的适用范围，并将研究者的政策提案与已实施政策分开记录。

\newpage
\section{探索性对照实验与结果分析}

\subsection{实验设置与主要结果}

为初步检查当前图谱对事件分析的作用，开展两轮小规模对照实验。两组均采用新会话，使用相同的请求配置\texttt{gpt-6.1-sol / high}；A组仅依据模型先验及公共题面，B组可额外检索冻结图谱。两组题面相同，均按任务要求禁止联网及读取原文、参考答案；图谱访问由只读工具记录。本次尚未建立操作系统级访问隔离。

\lead{第一轮：一般机制解释。}
从专业分析库分层随机抽取10篇文章，形成10道情境解释题。题目和参考要点在答题前锁定，答案匿名评分，两组均为100分。复查发现，题目省略了主体、时点和市场基线，同时给出较充分的机制线索，形成明显的得分天花板。

\lead{第二轮：具体事件与预期形成。}
另选10件事件，保留主体、信息时点、已有观测和判断期限，组成4道正向预期题、3道逆向解释题和3道选择题，资料截点为2026年11月4日。选择题答案与参考要点在答题前锁定，最终每题接受两份匿名评分。A组96.25分，B组95.75分；选择题均答对3题。逐题比较，B组1题较高、7题持平、2题较低。

\begin{center}
\small
\begin{tabularx}{\textwidth}{@{}YrrL{0.28\textwidth}@{}}
\toprule
测试 & 仅先验 & 先验＋图谱 & 主要观察 \\
\midrule
第一轮：10道机制题 & 100.00 & 100.00 & 存在明显天花板。 \\
第二轮：10道事件题 & 96.25 & 95.75 & 未观察到得分增益。 \\
\bottomrule
\end{tabularx}
\end{center}

第二轮B组调用32次检索或读取工具，阅读14个不同文件；17处引用均可定位且有读取记录。在CoreWeave题中，B组展开合同保障与融资条件，但遗漏题面给出的实际收费基准；AI CDS题中，两组均未充分利用跨企业差异约束解释。访问记录支持定位这类知识应用问题。

\subsection{解释边界与研究启示}

\lead{结果的适用范围。}
两轮实验均未观察到当前图谱的得分增益。样本量小、分数接近上限，第二轮0.5分差距不足以确立稳定效果。现有评分评价给定信息下的解释与预期形成质量；16条条件性预测已保存，结果尚未检验。以上结果不能用于报告市场方向命中率或投资收益。

\lead{后续需要分解检验的环节。}
知识积累的作用取决于相关内容是否存在、能否检索到，以及是否被正确应用。当前图谱以一般概念和机制为主，事件专门信息较多来自公共题面，因此需要进一步测试具体主体、时间联系和历史判断的积累是否提供有效信息。后续将增加原文检索对照，分别记录覆盖、检索、应用及阅读成本。

\newpage
\section{当前问题与后续研究安排}

\subsection{需要解决的主要问题}

\lead{知识整合与持续维护。}
三项新闻试点已形成解释稿，但尚需逐项确定主要维护位置，处理重复定义、对象继承和事件链接，并在接入后检查整体连贯性。下一阶段将先完成这三项整合，校准从文章到正式条目的完整流程，再扩展关联材料。新闻阅读仍承担基础概念与关系的建设职责。

\lead{可检验的研究问题。}
后续重点考察：当任务需要跨事件获取背景、还原此前判断并随新证据更新时，持续维护的金融认知图谱能否帮助模型形成更准确、完整且有依据的分析？系统文献比较将围绕背景选择、知识组织与证据更新展开。ReAct的信息获取过程\cite{react}及FinSearchComp的金融检索任务\cite{finsearch}为方法比较提供参考，具体实验仍需结合案例设计。

\lead{时间边界与评估有效性。}
限制历史任务的外部资料，仍不能保证模型先验不包含后续结果。因此，历史解释测试和前瞻预测检验将分别设计，固定资料截点、条件与观测期限，并保留未参与建图的测试材料。ForecastBench在结果未知时提交预测的设置\cite{forecastbench}可供参考。教师PPT测试需区分此前参与建图的内容与保留材料。

\subsection{后续约五周的工作计划}

以下安排对应开题十周计划的剩余时段。扩展规模根据整合质量与实验诊断调整。

\begin{center}
\small
\begin{tabularx}{\textwidth}{@{}L{0.16\textwidth}Y L{0.28\textwidth}@{}}
\toprule
时间 & 主要工作 & 预期产出 \\
\midrule
11月14—20日 & 完成三项试点的条目制作与整合审阅，统一知识归属及更新规则。 & 正式条目或更新；修订后的制作与审阅规范。 \\
11月21—27日 & 扩展少量关联事件，积累主体状态、机构分歧与判断修正；补充文献比较。 & 可追踪的时间线案例；研究问题与方法对照。 \\
11月28日—12月4日 & 冻结保留材料，比较仅模型先验、原文检索与图谱检索，设计分阶段信息更新任务。 & 预先确定的题目、评分规则、访问记录和对照结果。 \\
12月5—18日 & 分析失败案例、知识增量与检索成本，开展保留教师课件测试并完成总结。 & 评估与局限分析；文献综述及后续研究方案。 \\
\bottomrule
\end{tabularx}
\end{center}

对照实验将尽量对齐来源范围、信息截点、模型配置与检索预算，报告单题差异及失败原因。若扩展后的知识仍无可观察收益，将据此调整内容选择、组织方式或研究范围。阶段目标是形成可复核的案例与实验依据，为最终课题和方法选择提供支持。

\newpage
\begin{thebibliography}{9}
\small
\setlength{\itemsep}{7pt}
\setlength{\parskip}{0pt}

\bibitem{carlin}
Wendy Carlin and David Soskice.
\textit{Macroeconomics: Institutions, Instability, and Inequality}.
Oxford University Press, 2024. ISBN 9780198838661.
本项目采用所取得版本的导言及第1—18章，保留具体采用页码。

\bibitem{cecchetti}
Stephen G. Cecchetti and Kermit L. Schoenholtz.
\textit{Money, Banking and Financial Markets}.
6th edition, International Student Edition. McGraw-Hill Education, \copyright{}2021.
ISBN 9781260571363. 本项目实际采用第1—23章。

\bibitem{aicds}
鲍奕龙．2008再现！一文读懂华尔街最热门交易：AI CDS．华尔街见闻，页面显示日期2026年10月31日．
\href{https://wallstreetcn.com/articles/3784900}{公开文章原文}．

\bibitem{korea}
杨宸．韩国跌出“历史级便宜”，为何无人接盘？华尔街见闻，页面显示日期2026年10月20日．
\href{https://wallstreetcn.com/articles/3784053}{公开文章原文}．

\bibitem{basis}
张瑜．一文读懂隔夜美债崩盘的“罪魁”——基差交易．华尔街见闻转载并删减，页面显示日期2025年4月9日，页面署名李笑寅．原文为《对冲基金基差交易引发美债市场脆弱性担忧——海外周报第84期》．
\href{https://wallstreetcn.com/articles/3744871}{公开文章原文}．

\bibitem{react}
Shunyu Yao, Jeffrey Zhao, Dian Yu, et al.
ReAct: Synergizing Reasoning and Acting in Language Models.
\textit{International Conference on Learning Representations}, 2023.
\href{https://arxiv.org/abs/2210.03629v3}{arXiv:2210.03629v3}.

\bibitem{finsearch}
Liang Hu, Jianpeng Jiao, Jiashuo Liu, et al.
FinSearchComp: Towards a Realistic, Expert-Level Evaluation of Financial Search and Reasoning.
2025. \href{https://arxiv.org/abs/2509.13160v1}{arXiv:2509.13160v1}.

\bibitem{forecastbench}
Ezra Karger, Houtan Bastani, Chen Yueh-Han, et al.
ForecastBench: A Dynamic Benchmark of AI Forecasting Capabilities.
\textit{International Conference on Learning Representations}, 2025.
\href{https://arxiv.org/abs/2409.19839v5}{arXiv:2409.19839v5}.

\end{thebibliography}


\vspace{8pt}
{\small
\lead{报告依据。}
进展统计截至2026年11月13日。正式图谱节点与新闻试作清单分别计数。教材采用位置、材料原文与图表、试作成果、实验题面、答卷、访问记录及评分随项目保存，可供复核。新闻日期采用页面显示日期，机构观点按正文归属；文献出版与版本信息沿用2026年11月12日核对记录。
}

\end{document}
